\documentclass[12pt]{article}
\usepackage[english, bulgarian]{babel}
\usepackage[utf8]{inputenc}
\usepackage[T2A]{fontenc}
\usepackage{amssymb}
\usepackage{hyperref, fancyhdr, lastpage, fancyvrb, tcolorbox, titlesec}
\usepackage{array, tabularx, colortbl}
\usepackage{tikz}
\usepackage{venndiagram}
\usepackage{amsthm, bm}
\usepackage{relsize}
\usepackage{amsmath,physics}
\usepackage{mathtools}
\usepackage{subcaption}
\usepackage{scalerel}
\usepackage{theoremref}
\usepackage{circuitikz}
\usepackage{geometry}
\usepackage{stmaryrd}
\usepackage{forest}
\usepackage{cancel}
\usepackage{faktor}
\usepackage{gensymb}
\usetikzlibrary{automata, arrows, positioning, shapes}
\useforestlibrary{linguistics}

\ExplSyntaxOn
\NewDocumentCommand{\opair}{m}
 {
  \langle\mspace{2mu}
  \clist_set:Nn \l_tmpa_clist { #1 }
  \clist_use:Nn \l_tmpa_clist {,\mspace{3mu plus 1mu minus 1mu}\allowbreak}
  \mspace{2mu}\rangle
}
\ExplSyntaxOff

\hypersetup{
    colorlinks=true,
    linktoc=all,
    linkcolor=blue
}

\setlength\parindent{0pt}

\newcommand{\N}{\mathbb{N}}
\newcommand{\Z}{\mathbb{Z}}
\newcommand{\Q}{\mathbb{Q}}
\newcommand{\R}{\mathbb{R}}
\newcommand{\E}{\mathbb{E}}

\newcommand{\vars}{\operatorname{Vars}}
\newcommand{\free}{\operatorname{Free}}

\newcommand{\logand}{\; \& \;}

\newcommand{\calA}{\mathcal{A}}
\newcommand{\calS}{\mathcal{S}}
\newcommand{\calD}{\mathcal{D}}
\newcommand{\calL}{\mathcal{L}}
\newcommand{\calZ}{\mathcal{Z}}
\newcommand{\calQ}{\mathcal{Q}}
\newcommand{\calR}{\mathcal{R}}
\newcommand{\calN}{\mathcal{N}}
\newcommand{\calM}{\mathcal{M}}
\newcommand{\calF}{\mathcal{F}}
\newcommand{\calP}{\mathcal{P}}
\newcommand{\calC}{\mathcal{C}}
\newcommand{\calG}{\mathcal{G}}
\newcommand{\calB}{\mathcal{B}}

\newcommand{\dequiv}{\stackrel{\text{деф.}}{\longleftrightarrow}}

\newcommand{\db}[1]{\llbracket #1 \rrbracket}

\DeclareMathOperator*{\bigand}{\scalerel*{\&}{\prod}}
\DeclareMathOperator*{\bigor}{\scalerel*{\lor}{\sum}}

\newtheorem*{definition}{Дефиниция}
\newtheorem{problem}{Задача}[section]
\newtheorem*{claim}{Твърдение}
\newtheorem*{property}{Свойство}
\newtheorem*{hint}{Упътване}
\theoremstyle{definition}
\newtheorem*{solution}{Решение}
\newtheorem*{notation}{Нотация}

\title{(Не)определимост и изпълнимост -- 2025/2026}
\author{Тодор Дуков}
\date{}

\begin{document}
\maketitle

\section{(Не)определимост}

\begin{problem}
Дефинираме индуктивно $V_n$, за всяко естествено $n$:
\begin{align*}
	 & V_0        = \varnothing  \\
	 & V_{n + 1}  = \calP (V_n).
\end{align*}
След това полагаме:
\[
	V_\omega = \bigcup\limits_{n \in \N} V_n.
\]
Разглеждаме предикатния език $\calL$, съставен от двуместен предикатен символ $\operatorname{in}$, заедно с $\calL$-структурата $\calS = \opair{V_\omega, \operatorname{in}^\calS}$, където:
\[
	\operatorname{in}^\calS(X, Y) \dequiv X \in Y.
\]
Да се докаже, че в $\calS$ са определими:
\begin{itemize}
	\item графиките на операциите чифт, наредена двойка (дефиницията на Kuratowski) и степенно множество;
	\item графиките на операциите обединение и сечение;
	\item графиките на операциите събирание, умножение и степенуване;
	\item множеството от всички релации, множеството от всички функции и множеството от всички двойки равномощни множества;
	\item за всяко естествено $n$, множествата $\{ V_n \}, \{ n \}$ и множеството $\omega$;
\end{itemize}
Да се докаже, че ако една за релация $R \in V_n$, множеството $\{ R \}$ е определимо в $\calS$, то тогава $R$ е определимо в $\calS$.
\end{problem}

\begin{problem}
Нека фиксираме някаква крайна азбука $\Sigma$.
Нека $\calL$ е предикатен език, съставен от двуместния предикатен символ $<$, заедно с едноместните предикатни символи $q_\sigma$, за всяко $\sigma \in \Sigma$.

За всяко $\alpha \in \Sigma^*$, дефинираме $\calL$-структурата $\calA_\alpha$, за която:
\begin{align*}
	|\calA_\alpha|             & = \{ 0, \dots, |\alpha| \}                          \\
	i <^{\calA_\alpha} j       & \dequiv i < j                                       \\
	q_\sigma^{\calA_\alpha}(i) & \dequiv i < |\alpha| \text{ и } \alpha[i] = \sigma.
\end{align*}
За една формула $\varphi$ дефинираме:
\[
	\operatorname{Lang}[\varphi] = \{ \alpha \in \Sigma^* \mid \calA_\alpha \models \varphi \}.
\]
Да се докаже, че ако един език $L$ е регулярен и в конструкцията му не участва звезда на Клини, то тогава има формула $\varphi$, за която $\operatorname{Lang}[\varphi] = L$.
\end{problem}

\begin{problem}
Нека $\Sigma = \{ a, b \}$.
Разглеждаме предикатния език $\calL$, съставен от триместния предикатен символ $\operatorname{cat}$ и едноместния предикатен символ $\operatorname{leq}$, заедно със $\calL$-структурата $\calM = \opair{\Sigma^*, \operatorname{cat}^\calM, \operatorname{leq}^\calM}$, където:
\begin{align*}
	\operatorname{cat}(\alpha, \beta, \gamma) & \dequiv \alpha \cdot \beta = \gamma \\
	\operatorname{leq}(\alpha)                & \dequiv |\alpha|_a \leq |\alpha|_b.
\end{align*}
Да се докаже, че в $\calM$ са определими:
\begin{itemize}
	\item за всяко $\alpha \in \Sigma^*$, множеството $\{ \alpha \}$;
	\item езика $\{ ab \}^*$;
	\item езика $\{ a \}^* \cdot \{ b \}^*$;
	\item езика $\{ a^n b^n \mid n \in \N \}$;
	\item езика $\{ a^n b^n a^n \mid n \in \N \}$.
\end{itemize}
Да се определят всички автоморфизми в $\calM$.
\end{problem}

\newpage 

\begin{problem}
Нека $\Sigma$ е крайна азбука.
Със $\operatorname{PC}(\Sigma)$ да бележим множеството от всички езици $T \subseteq \Sigma^*$, които се затворени надолу отностно $\preceq_{\operatorname{pref}}$.
Разглеждаме предикатния език $\calL$, съставен от двуместния предикатен символ $\operatorname{sub}$, заедно със $\calL$-структурата $\calM_\Sigma = \opair{\operatorname{PC}(\Sigma), \operatorname{sub}^{\calM_\Sigma}}$, където:
\[
  \operatorname{sub}^{\calM_\Sigma}(T_1, T_2) \dequiv T_1 \subseteq T_2.
\]
Да се докаже, че в $\calM_\Sigma$ са определими:
\begin{itemize}
  \item $\{ \Sigma^* \}$;
  \item за всяко $n$, множеството $\{ \Sigma^{\leq n} \}$;
  \item множеството от тези езици, които са вериги в $\opair{\Sigma^*, \preceq_{\operatorname{pref}}}$;
  \item множеството от тези безкрайни езици, които са вериги в $\opair{\Sigma^*, \preceq_{\operatorname{pref}}}$.
\end{itemize}
Да се докаже, че за всяко естествено $n \geq 1$ има формула $\varphi_n$, за която:
\[
  \calM_\Sigma \models \varphi_n \longleftrightarrow |\Sigma| = n.
\]
Да се определят всички автоморфизми в $\calM_\Sigma$.
\end{problem}

\newpage

\section{Изпълнимост}

\begin{problem}
Да се докаже, че следното множество от формули е изпълнимо:
\begin{align*}
	 & \varphi_1 : \forall x \forall y \exists ! z \exists ! t (f(x, z) \doteq y \logand f(t, x) \doteq y) \\
	 & \varphi_2 : \neg \forall x \forall y \forall z (f(x, f(y, z)) \doteq f(f(x, y), z))                 \\
	 & \varphi_3 : \forall x (f(x, a) \doteq x \logand f(a, x) \doteq x)                                   \\
	 & \varphi_4 : \forall x \forall y (f(x, y) \doteq f(y, x)).
\end{align*}
\end{problem}

\begin{problem}
Да се докаже, че следното множество от формули е изпълнимо:
\begin{align*}
	 & \varphi_1 : \forall x \forall y \forall z (f(x, f(y, z)) \doteq f(f(x, y), z))       \\
	 & \varphi_2 : \forall x \forall y \forall z (g(x, g(y, z)) \doteq g(g(x, y), z))       \\
	 & \varphi_3 : \forall x \forall y \forall z (g(x, f(y, z)) \doteq f(g(x, y), g(x, z))) \\
	 & \varphi_4 : \forall x \forall y (f(x, g(y, x)) \doteq x)                             \\
	 & \varphi_5 : \exists x \exists y (f(x, y) \not \doteq g(x, y))                        \\
\end{align*}
\end{problem}

\begin{problem}
Да се докаже, че следното множество от формули е изпълнимо:
\begin{align*}
	 & \varphi_1 : \forall x \forall y \forall z (f(x, f(y, z)) \doteq f(f(x, y), z))       \\
	 & \varphi_2 : \forall x \forall y \forall z (g(x, g(y, z)) \doteq g(g(x, y), z))       \\
	 & \varphi_3 : \forall x \forall y \forall z (g(x, f(y, z)) \doteq f(g(x, y), g(x, z))) \\
	 & \varphi_4 : \exists x \exists y (p(x) \logand p(y) \logand p(f(x, y)))               \\
	 & \varphi_5 : \exists x \exists y (p(x) \logand p(y) \logand \neg p(f(x, y)))          \\
	 & \varphi_6 : \exists x \exists y (p(x) \logand \neg p(y) \logand p(f(x, y)))          \\
	 & \varphi_7 : \exists x \exists y (p(x) \logand \neg p(y) \logand \neg p(f(x, y)))     \\
	 & \varphi_8 : \forall x \forall y (p(x) \logand p(y) \Rightarrow \neg p(g(x, y)))
\end{align*}
\end{problem}

\begin{problem}
Да се докаже, че следното множество от формули е изпълнимо:
\begin{align*}
	 & \varphi_1 : \forall x \neg p(x, x)                                                     \\
   & \varphi_2 : \forall x \forall y \forall z (p(x, y) \logand p(y, z) \Rightarrow p(x, z)) \\
   & \varphi_3 : \forall x \forall (q(x, y) \Leftrightarrow p(x, y) \logand \neg \exists z (p(x, z) \logand p(z, y))) \\
   & \psi_n : \exists x \exists y_1 \dots \exists y_n (\bigand_{1 \leq i \leq n} q(x, y_i) \logand \bigand_{1 \leq i, j \leq n} \neg q(y_i, y_j) \logand \forall y (q(x, y) \Rightarrow \bigor_{1 \leq i \leq n} y \doteq y_i))
\end{align*}
\end{problem}

\end{document}
